\CWHeader{Лабораторная работа \textnumero 7}

\CWProblem{
Имеется натуральное число $n$. За один ход с ним можно произвести следующие действия:

\begin{itemize}
    \item вычесть единицу;
    \item разделить на два;
    \item разделить на три.
\end{itemize}

При этом стоимость каждой операции~--- текущее значение $n$. Стоимость преобразования~--- суммарная стоимость всех операций в преобразовании.
Вам необходимо с помощью последовательностей указанных операций преобразовать число $n$ в единицу таким образом, чтобы стоимость
преобразования была наименьшей. Делить можно только нацело.

\textbf{Формат ввода}

В первой строке задано $2 \le n \le 10^7$.

\textbf{Формат вывода}

Выведите на первой строке искомую наименьшую стоимость. Во второй строке должна содержаться последовательность операций.
Если было произведено деление на 2 или на 3, выведите \texttt{/2} (или \texttt{/3}). Если же было вычитание, выведите \texttt{-1}.
Все операции выводите, разделяя пробелом.
}
\pagebreak